
\documentclass{article}
\usepackage{colortbl}
\usepackage{makecell}
\usepackage{multirow}
\usepackage{supertabular}

\begin{document}

\newcounter{utterance}

\centering \large Interaction Transcript for game `ifeval', experiment `ifeval', episode 436 with reflexion{-}llama{-}3.1{-}8b{-}instruct.
\vspace{24pt}

{ \footnotesize  \setcounter{utterance}{1}
\setlength{\tabcolsep}{0pt}
\begin{supertabular}{c@{$\;$}|p{.15\linewidth}@{}p{.15\linewidth}p{.15\linewidth}p{.15\linewidth}p{.15\linewidth}p{.15\linewidth}}
    \# & \multicolumn{2}{c}{Player} && \multicolumn{2}{c}{Game Master} \\
    \hline

    \theutterance \stepcounter{utterance}  
    & & & \multicolumn{4}{p{0.6\linewidth}}{
        \cellcolor[rgb]{0.9,0.9,0.9}{
            \makecell[{{p{\linewidth}}}]{
                \texttt{\tiny{[P1$\langle$GM]}}
                \texttt{You have attempted this same task before and did not succeed. These are the plans you wrote after those attempts.} \\
\texttt{Trial \#1: मेरी पिछली कोशिश में मैंने वित्तीय संकट के वर्तमान संकट के पर्यावरण पर प्रभाव के बारे में एक निबंध लिखने का प्रयास किया, लेकिन मैंने अपने लक्ष्य को पूरा नहीं किया। मेरी गलती यह थी कि मैंने अपने निबंध में "मौसम", "ऊर्जा", और "हरा" शब्दों को शामिल नहीं किया।} \\
\\ 
\texttt{नई योजना: अगली बार, मैं अपने निबंध में "मौसम", "ऊर्जा", और "हरा" शब्दों को शामिल करने के लिए विशेष ध्यान दूंगा। मैं अपने निबंध को विस्तार से लिखूंगा और वित्तीय संकट के वर्तमान संकट के पर्यावरण पर प्रभाव के बारे में विस्तार से चर्चा करूंगा। मैं अपने निबंध में नवीकरणीय ऊर्जा स्रोतों के महत्व, ऊर्जा की बचत के तरीके, और हरित प्रौद्योगिकी के उपयोग के बारे में भी चर्चा करूंगा।} \\
\\ 
\\ 
\texttt{First repeat the request below word for word without change, then give your answer. Do not say any words or characters before repeating the request.} \\
\\ 
\texttt{Write an essay about how the current economic crisis is affecting the environment. In your essay, include the keywords: "climate", "energy", and "green". Make sure your entire response is in Hindi, no other language is allowed.} \\
            }
        }
    }
     \\ \\

    \theutterance \stepcounter{utterance}  
    & \multicolumn{4}{p{0.6\linewidth}}{
        \cellcolor[rgb]{0.561,0.737,0.561}{
            \makecell[{{p{\linewidth}}}]{
                \texttt{\tiny{[P1$\rangle$GM]}}
                \texttt{मैं अपने निबंध में "मौसम", "ऊर्जा", और "हरा" शब्दों को शामिल करने के लिए विशेष ध्यान दूंगा। मैं अपने निबंध को विस्तार से लिखूंगा और वित्तीय संकट के वर्तमान संकट के पर्यावरण पर प्रभाव के बारे में विस्तार से चर्चा करूंगा। मैं अपने निबंध में नवीकरणीय ऊर्जा स्रोतों के महत्व, ऊर्जा की बचत के तरीके, और हरित प्रौद्योगिकी के उपयोग के बारे में भी चर्चा करूंगा।} \\
\\ 
\texttt{वित्तीय संकट के वर्तमान संकट ने पर्यावरण पर गहरा प्रभाव डाला है। मौसम के परिवर्तन के कारण जलवायु परिवर्तन की समस्या बढ़ रही है, जिससे पृथ्वी के तापमान में वृद्धि हो रही है। ऊर्जा की मांग बढ़ने से ऊर्जा के नवीकरणीय स्रोतों की आवश्यकता बढ़ रही है, जैसे कि सौर ऊर्जा, पवन ऊर्जा, और जलविद्युत ऊर्जा। इन नवीकरणीय ऊर्जा स्रोतों का उपयोग करके हम ऊर्जा की बचत कर सकते हैं और पर्यावरण को नुकसान से बचा सकते हैं।} \\
\\ 
\texttt{हरित प्रौद्योगिकी का उपयोग करके हम पर्यावरण को संरक्षित कर सकते हैं। हरित प्रौद्योगिकी के माध्यम से हम ऊर्जा की बचत कर सकते हैं, जलवायु परिवर्तन को रोक सकते हैं, और पर्यावरण को स्वच्छ बना सकते हैं। इसके अलावा, हरित प्रौद्योगिकी का उपयोग करके हम अपने जीवनशैली को भी बदल सकते हैं और पर्यावरण के प्रति जागरूकता बढ़ा सकते हैं।} \\
\\ 
\texttt{नवीकरणीय ऊर्जा स्रोतों का उपयोग करके हम ऊर्जा की बचत कर सकते हैं और पर्यावरण को नुकसान से बचा सकते हैं। सौर ऊर्जा का उपयोग करके हम अपने घरों में ऊर्जा की बचत कर सकते हैं और पर्यावरण को स्वच्छ बना सकते हैं। पवन ऊर्जा का उपयोग करके हम ऊर्जा की बचत कर सकते हैं और पर्यावरण को नुकसान से बचा सकते हैं।} \\
\\ 
\texttt{अंत में, वित्तीय संकट के वर्तमान संकट ने पर्यावरण पर गहरा प्रभाव डाला है। मौसम के परिवर्तन के कारण जलवायु परिवर्तन की समस्या बढ़ रही है, जिससे पृथ्वी के तापमान में वृद्धि हो रही है। ऊर्जा की मांग बढ़ने से ऊर्जा के नवीकरणीय स्रोतों की आवश्यकता बढ़ रही है। इसलिए, हमें नवीकरणीय ऊर्जा स्रोतों का उपयोग करना चाहिए, ऊर्जा की बचत करनी चाहिए, और हरित प्रौद्योगिकी का उपयोग करना चाहिए ताकि हम पर्यावरण को संरक्षित कर सकें और जलवायु परिवर्तन को रोक सकें।} \\
            }
        }
    }
    & & \\ \\

    \theutterance \stepcounter{utterance}  
    & & & \multicolumn{2}{p{0.3\linewidth}}{
        \cellcolor[rgb]{0.9,0.9,0.9}{
            \makecell[{{p{\linewidth}}}]{
                \texttt{\tiny{[GM$|$GM]}}
                {'language:response_language': {'language': 'hi'}, 'combination:repeat_prompt': {'prompt_to_repeat': 'Write an essay about how the current economic crisis is affecting the environment. In your essay, include the keywords: "climate", "energy", and "green". Make sure your entire response is in Hindi, no other language is allowed.'}, 'keywords:existence': {'keywords': ['climate', 'energy', 'green']}}
            }
        }
    }
    & & \\ \\

    \theutterance \stepcounter{utterance}  
    & & & \multicolumn{2}{p{0.3\linewidth}}{
        \cellcolor[rgb]{0.9,0.9,0.9}{
            \makecell[{{p{\linewidth}}}]{
                \texttt{\tiny{[GM$|$GM]}}
                \texttt{game\_result = LOSE} \\
            }
        }
    }
    & & \\ \\

    \theutterance \stepcounter{utterance}  
    & & \multicolumn{4}{p{0.6\linewidth}}{
        \cellcolor[rgb]{0.95,0.95,0.95}{
            \makecell[{{p{\linewidth}}}]{
                \texttt{\tiny{[P1$\leftrightarrow$P]}}
                \texttt{Post{-}episode Reflexion (not sent to GM):} \\
\texttt{नई योजना: मेरी पिछली कोशिश में मैंने वित्तीय संकट के वर्तमान संकट के पर्यावरण पर प्रभाव के बारे में एक निबंध लिखने का प्रयास किया, लेकिन मैंने अपने लक्ष्य को पूरा नहीं किया। मेरी गलती यह थी कि मैंने अपने निबंध में "मौसम", "ऊर्जा", और "हरा" शब्दों को शामिल नहीं किया और मैंने अपने निबंध को विस्तार से लिखने के लिए पर्याप्त समय नहीं दिया।} \\
\\ 
\texttt{अगली बार, मैं अपने निबंध में "मौसम", "ऊर्जा", और "हरा" शब्दों को शामिल करने के लिए विशेष ध्यान दूंगा। मैं अपने निबंध को विस्तार से लिखूंगा और वित्तीय संकट के वर्तमान संकट के पर्यावरण पर प्रभाव के बारे में विस्तार से चर्चा करूंगा। मैं अपने निबंध में नवीकरणीय ऊर्जा स्रोतों के महत्व, ऊर्जा की बचत के तरीके, और हरित प्रौद्योगिकी के उपयोग के बारे में भी चर्चा करूंगा।} \\
\\ 
\texttt{इसके अलावा, मैं अपने निबंध को विभिन्न भागों में विभाजित करूंगा, जैसे कि परिचय, मुख्य भाग, और निष्कर्ष। मैं अपने निबंध को विस्तार से लिखने के लिए पर्याप्त समय दूंगा और अपने विचारों को स्पष्ट और संक्षिप्त तरीके से व्यक्त करूंगा।} \\
\\ 
\texttt{अंत में, मैं अपने निबंध को एक संगठित और स्पष्ट तरीके से लिखूंगा, जिससे पाठकों को आसानी से समझने में मदद मिलेगी।} \\
            }
        }
    }
    & \\ \\

\end{supertabular}
}

\end{document}
